\chapter{A team analysis: the mystery resonance}
\label{ch:team}

\chaptergoals{6 of 6 in Part~I: everything together.}
{\cref{ch:git,ch:branches,ch:remotes,ch:pixi,ch:ci}; a GitHub account; some
Python. Still no SHiP software: the team repository brings its own pixi
environment with PyROOT.}
{take one piece of a shared analysis from an empty function to a merged pull
request, through local tests, CI and review, while the team's analysis grows
with every merge.}

The whole group writes one small analysis together. Each person writes one
piece, in a pull request that CI checks and a teammate reviews, and the
analysis only works when everybody's pieces fit. The exercise runs in rounds;
each round uses one of the earlier chapters.

\begin{center}\small
\begin{tabular}{l l l}
\toprule
Round & What you do & Uses \\
\midrule
0 & fork, clone, install the environment & \cref{ch:remotes,ch:pixi} \\
1 & write and test your task on your machine & \cref{ch:git,ch:pixi} \\
2 & open the pull request, get the checks green & \cref{ch:remotes,ch:ci} \\
3 & review, rebase, merge & \cref{ch:branches,ch:remotes} \\
4 & watch the analysis grow & \cref{ch:remotes,ch:pixi} \\
5 & find the particle & PyROOT \\
\bottomrule
\end{tabular}
\end{center}

\section{The data}

The file \file{mystery.root} holds 20\,000 measurements of the invariant mass
of the decay products of an unknown particle, in GeV. They are stored in an
RNTuple called \texttt{events}, with a single field, \texttt{mass}, of type
\texttt{float}. The line shape is a relativistic resonance, described by a
Breit--Wigner distribution:
\[
  f(m) \;=\; \frac{1}{\pi}\,\frac{\Gamma/2}{(m-M)^2 + \Gamma^2/4},
\]
with mass $M$ and width $\Gamma$ (the full width at half maximum). Your team
has to measure $M$ and $\Gamma$ with their uncertainties, make a plot of the
fit, and find out which particle this is. Nobody is told the answer.

\section{One script, ten tasks}

The analysis lives in one file, \file{mystery\_analysis.py}, in a repository
of its own, the \emph{team repository} (\cref{sec:teamsetup}). It is cut into ten tasks, each one a Python
function that is empty to start with:

\begin{center}\small
\begin{tabularx}{\linewidth}{r l L}
\toprule
& Function & Input $\to$ Output \\
\midrule
1  & \texttt{open\_file}         & file name $\to$ RNTuple reader \\
2  & \texttt{read\_masses}       & reader $\to$ list of masses \\
3  & \texttt{summarise}          & masses $\to$ count, mean, median, central 68\% interval \\
4  & \texttt{book\_histogram}    & summary $\to$ empty histogram \texttt{h\_mystery} \\
5  & \texttt{fill\_histogram}    & histogram, masses $\to$ filled histogram \\
6  & \texttt{make\_model}        & histogram $\to$ Breit--Wigner \texttt{f\_func} with starting values \\
7  & \texttt{fit\_model}         & histogram, model $\to$ fit result \\
8  & \texttt{extract\_results}   & fitted model $\to$ mass, width, errors, fit quality \\
9  & \texttt{identify\_particle} & results $\to$ best-matching known particle \\
10 & \texttt{save\_plot}         & histogram, model, results $\to$ \file{mystery\_fit.pdf} \\
\bottomrule
\end{tabularx}
\end{center}

Above each function, a comment gives the contract: what the function
receives (\emph{Input}), what it must return (\emph{Output}), the ROOT calls
that are useful (\emph{Relevant}), some \emph{Hints}, and how to test it. It
tells you what to use, not what to write. Task~7, for example:

\codefile[linerange={193-214},basicstyle=\ttfamily\footnotesize]{Python}{code/teamwork/template/mystery_analysis.py}

The contract is what makes the pieces fit: the owner of task~8 relies on task~7
returning exactly what its Output line promises. Do not rename your function or
change its arguments, and keep your code between the two marker lines.

\subsection*{Running the chain}

At the bottom of the file, a function written by the organisers calls the ten
tasks in order. It stops, without an error, at the first task that is not
written yet:

\begin{outputlisting}
$ pixi cvmfs analysis
Pixi task (analysis): python mystery_analysis.py
--- task  1/10: open_file
mystery.root: 20000 entries in 'events'
--- task  2/10: read_masses
--- task  3/10: summarise
  n = 2e+04, min = 20.03, max = 199.7, mean = ..., median = ...
--- task  4/10: book_histogram
--- task  5/10: fill_histogram
--- task  6/10: make_model is not written yet; the chain stops here (5 of 10 tasks done)
\end{outputlisting}

Every merged pull request lets the chain run one step further, until it ends
with \texttt{all 10 tasks done}.

\subsection*{Testing your task on its own}

You do not have to wait for your teammates. \file{check\_task.py} tests one
task with inputs it makes up itself, from a toy Breit--Wigner sample whose
answer it knows ($M = 50$\,GeV, $\Gamma = 4$\,GeV: \emph{not} the mystery
particle), and checks your function's Output against its contract:

\begin{outputlisting}
$ pixi cvmfs check-task 4
Pixi task (check-task): python check_task.py 4
task  4: FAILED
      the histogram must be named 'h_mystery', got 'h'
$ pixi cvmfs check-task 4
Pixi task (check-task): python check_task.py 4
task  4: passed
\end{outputlisting}

The same checker, for all ten tasks at once, is what CI runs on every pull
request (\cref{sec:teamci}).

\section{Who does what}
\label{sec:teamroles}

With ten participants, everyone takes one task. With fewer, some take two or
three neighbouring tasks; with more, pairs share a task, or the extra people
take the bonus tasks of \cref{sec:bonus}.

\begin{center}\small
\begin{tabular}{c l}
\toprule
People & Tasks per person \\
\midrule
10 & 1 \,|\, 2 \,|\, 3 \,|\, 4 \,|\, 5 \,|\, 6 \,|\, 7 \,|\, 8 \,|\, 9 \,|\, 10 \\
9  & 1 \,|\, 2 \,|\, 3 \,|\, 4 \,|\, 5 \,|\, 6 \,|\, 7 \,|\, 8+9 \,|\, 10 \\
8  & 1 \,|\, 2+3 \,|\, 4 \,|\, 5 \,|\, 6 \,|\, 7 \,|\, 8+9 \,|\, 10 \\
7  & 1 \,|\, 2+3 \,|\, 4+5 \,|\, 6 \,|\, 7 \,|\, 8+9 \,|\, 10 \\
6  & 1+2+3 \,|\, 4+5 \,|\, 6 \,|\, 7 \,|\, 8+9 \,|\, 10 \\
5  & 1+2+3 \,|\, 4+5 \,|\, 6 \,|\, 7+8 \,|\, 9+10 \\
\bottomrule
\end{tabular}
\end{center}

\textbf{Your reviewer is the owner of the next task}, the person who will use
what your function returns (task~10 is reviewed by the owner of task~1).
Write your name and your reviewer's next to your task in the table of the
repository's \file{README.md}, in your first commit.

\Needspace{12\baselineskip}
\section{Round 0: the repository and its environment}
\label{sec:teamworkflow}

Fork the team repository on GitHub (``Fork'' button), then:
\begin{shell}
git clone git@github.com:<you>/sw_tutorial_2026.git
cd sw_tutorial_2026
git remote add upstream git@github.com:matclim/sw_tutorial_2026.git
cp <wherever-it-is>/mystery.root .          # data: not under version control
pixi cvmfs analysis                         # stops at task 1 for now
pixi cvmfs check                            # 10 not written yet, 0 failed
\end{shell}
\cmd{pixi task list} shows what else the repository offers.

\begin{checkpoint}
\cmd{pixi cvmfs analysis} stops at task~1 without an error, and
\cmd{pixi cvmfs check} reports \texttt{0 passed, 10 not written yet, 0 failed}.
\end{checkpoint}

\Needspace{12\baselineskip}
\section{Round 1: your task, on your machine}

One branch, one task:
\begin{shell}
git switch -c task-04-book-histogram        # task-NN-short-description
vim mystery_analysis.py                     # write your function
vim README.md                               # your name and your reviewer's
pixi cvmfs check-task 4                     # until it says "passed"
pixi cvmfs lint                             # until ruff is happy
git add mystery_analysis.py README.md
git commit -m "feat(task-04): book the mystery histogram"
\end{shell}
Commit as often as you like on your branch: small steps are easier to review
and to undo.

\Needspace{12\baselineskip}
\section{Round 2: the pull request and its checks}

\begin{shell}
git push -u origin task-04-book-histogram
\end{shell}
Open the pull request against \texttt{main} of the team repository (the page
GitHub shows after the push has a button for it). Describe what you did,
paste the output of \cmd{pixi cvmfs check-task 4}, and ask your reviewer for a
review.

CI starts by itself (\cref{sec:teamci}). Wait for ``All checks passed''. If it
is red, open the log of the failing job, reproduce the failure with the same
pixi task on your machine, fix it, commit, and \cmd{git push}: the pull
request updates and CI runs again.

\Needspace{12\baselineskip}
\section{Round 3: review, rebase, merge}

\paragraph{As a reviewer.} Read the change in the ``Files changed'' tab. Does
it do what the contract says, and nothing else? Would you, as the next task's
owner, be happy to receive its output? Comment on lines you do not
understand. Approve only when you are satisfied and the checks are green.

\paragraph{As an author.} Answer every comment, change what needs changing,
and push again; a new push dismisses earlier approvals, so ask for a new one.

\paragraph{When your pull request is out of date.} Every merge moves
\texttt{main} forward, and GitHub then asks for your branch to be brought up
to date. Rebase:
\begin{shell}
git fetch upstream
git rebase upstream/main
pixi cvmfs check-task 4
git push --force-with-lease
\end{shell}
Since every task lives in its own function, this rebase almost never
conflicts. If it does, \cref{sec:mergerebase} explains what to do. An
approved, green, up-to-date pull request is merged by the organisers (or by
the merge queue, if the repository has one).

\Needspace{12\baselineskip}
\section{Round 4: watch the analysis grow}

After every merge, everybody can run the team's analysis as far as it goes:
\begin{shell}
git switch main && git pull upstream main
pixi cvmfs analysis
git switch task-04-book-histogram           # back to your own work
\end{shell}

\section{Round 5: the answer}

\begin{exercise}{Find the mystery particle}
When the chain reaches \texttt{all 10 tasks done}, open
\file{mystery\_fit.pdf}.
\begin{enumerate}
  \item What are the mass and the width, with their uncertainties?
  \item Which particle is it? Is the fitted width compatible with that
        particle's, too, or only the mass?
  \item Is the fit good? Look at $\chi^2/\mathrm{ndf}$ and at the tails on a
        logarithmic scale.
  \item Change one thing at a time (the binning in task~4, chi-square versus
        likelihood in task~7, the starting values in task~6) and rerun. Which
        changes move the result by more than its uncertainty?
\end{enumerate}
\end{exercise}

\section{Bonus tasks}
\label{sec:bonus}

For those who finish early, or for an eleventh and twelfth participant. Each
is a separate pull request.

\begin{description}[style=nextline]
\item[B1: pulls] Under the plot, add a second pad with
  $(n_i - f(m_i))/\sigma_i$ per bin. Does the model describe the tails?
\item[B2: chi-square versus likelihood] Fit with and without the \texttt{"L"}
  option and tabulate both results. Which is biased, and why, for bins with
  few entries?
\item[B3: a background] The organisers have a second sample,
  \file{mystery\_bkg.root}, with a falling background under the peak. Add an
  exponential to the model (\texttt{"[0]*TMath::BreitWigner(x,[1],[2]) +
  [3]*exp(-[4]*x)"}) and compare the result with and without it.
\item[B4: detector resolution] Real detectors smear the mass. The sample
  \file{mystery\_res.root} has a Gaussian resolution added; fit it with a
  Voigt profile (\texttt{TMath::Voigt}), the convolution of a Breit--Wigner
  with a Gaussian. What happens to the width if you ignore the resolution?
\end{description}

\section{Setting it up}
\label{sec:teamsetup}

The team repository is
\href{https://github.com/matclim/sw_tutorial_2026}{\file{github.com/matclim/sw\_tutorial\_2026}}.
It starts from the
files in \file{code/teamwork/template/}: the skeleton
\file{mystery\_analysis.py}, the checker \file{check\_task.py}, the pixi
manifest \file{pixi.toml} with its lock file, the lint settings
\file{ruff.toml}, the CI workflow \file{.github/workflows/ci.yml}, a
\file{README.md} with the task table, the rules in \file{CONTRIBUTING.md}, and
a \file{.gitignore} that keeps data, plots and \file{.pixi/} out of the
history.

\iforganiser\input{code/teamwork/organiser/team-organiser}\fi
